\documentclass[10pt,a4paper]{amsart}
\usepackage[margin=19mm]{geometry}
\usepackage{amsmath,amssymb,mathtools}
\usepackage[scr=rsfs,cal=euler]{mathalfa}
\usepackage{xfrac}
\usepackage[unicode,hidelinks,bookmarks=false]{hyperref}
\newcommand{\ie}{i.e.\ }
\newcommand{\cf}{cf.\ }
\newcommand{\resp}{respectively\ }
\newcommand{\Subsection}{\subsection}
\providecommand{\nobreakdash}{\nobreak\hbox{-}}
\setlength{\emergencystretch}{2em}
\multlinegap0pt
\usepackage{amssymb}
\usepackage[normalem]{ulem}
\usepackage{xcolor}
\usepackage{float}
\newtheorem{theorem}{Theorem}
\title{A simple second-order nonstandard numerical scheme for a two-stage structured species with recruitment}
\author{Manh Tuan Hoang}
\date{Source: arXiv:2608.09141v1 (CC BY 4.0)}
\begin{document}
\maketitle

\noindent\emph{Source and licence.} A simple second-order nonstandard numerical scheme for a two-stage structured species with recruitment. Version arXiv:2608.09141v1 (CC BY 4.0), distributed by arXiv under the Creative Commons Attribution 4.0 International licence, http://creativecommons.org/licenses/by/4.0/, as recorded in the arXiv licence metadata for this exact version and shown on the abstract page https://arxiv.org/abs/2608.09141v1. Full licence notice: \texttt{licenses/arxiv-2608.09141-CC-BY-4.0.txt}.\par\smallskip

\noindent\textit{Editorial scope.} Counted unit: the article's theorem Theorem1 (source0.tex lines 614--616). The article's complete original proof of it is delivered with it (source0.tex lines 619--645, one \texttt{proof} environment of 27 lines). No mathematics is written, repaired or re-derived: the statement and the proof are the article's own lines, in the article's own notation, and no retained line is substituted or re-flowed.  Retained uncounted as the source-local context the proof needs: lines 548--550; lines 585--590; lines 602--604.  Nothing was omitted from any retained span, so the package carries no omission record.  No retained line contains a citation, so the wrapper carries no bibliography and no bibliography entry.  No retained line contains a figure environment, an image-inclusion command or drawing code, so no image is delivered and the package declares no asset dependency.  Editorial formatting only, in addition: an AMS \texttt{amsart} preamble that restates the article's own declarations and macros used by the retained lines, namely the environments \texttt{theorem} on separate counters, together with the standard packages those lines need.  The retained environments print in this document as Theorem 1 in order of appearance, whereas the article prints the same items with the numbering of its own full document.  The complete native source of the article as distributed in the arXiv e-print for this exact version is bundled unchanged as \texttt{sources/207218/source0.tex}.
\begin{equation}\label{eq:3a}
\Omega = \{(x, y) \in \mathbb{R}^2|x, y \geq 0\}
\end{equation}

\begin{equation}\label{eq:NSFD}
\begin{split}
\dfrac{x_{k+1}- x_k}{\phi_1(h, x_k, y_k)} &= \delta y_k - \dfrac{\alpha x_{k}}{\beta + x_k}- \mu x_{k} + \tau_1x_k - \tau_1x_{k+1}, \\
\dfrac{y_{k+1} - y_k}{\phi_2(h, x_k, y_k)} &= \dfrac{\alpha x_{k}}{\beta + x_k} - (\mu+F) y_{k} + \tau_2y_k - \tau_2y_{k+1},
\end{split}
\end{equation}

\begin{equation}\label{eq:6}
\tau_1 \geq \dfrac{\alpha}{\beta} + \mu, \quad \tau_2 \geq \mu + F.
\end{equation}

\begin{theorem}[Positively invariant set]\label{Theorem1}
Under the condition \eqref{eq:6}, the discrete-time model \eqref{eq:NSFD} admits the set $\Omega$ defined in \eqref{eq:3a} as a positively invariant set.
\end{theorem}

\begin{proof}
Assume that $x_0, y_0 \geq$. We need to show that $x_k, y_k \geq$ for all $k > 0$. Indeed, \eqref{eq:NSFD} can be transformed into the form
\begin{equation}\label{eq:7}
\begin{split}
x_{k + 1} &= \dfrac{x_k + \phi_1\delta y_k - \phi_1\dfrac{\alpha x_{k}}{\beta + x_k}- \phi_1\mu x_{k} + \phi_1\tau_1x_k}{1 + \tau_1\phi_1} = \dfrac{x_k + \phi_1\delta y_k - \phi_1\bigg(\dfrac{\alpha}{\beta + x_k} +  \mu - \tau_1\bigg)}{1 + \tau_1\phi_1}\\
%%
%%
%%
y_{k + 1} &= \dfrac{y_k + \phi_2\dfrac{\alpha x_{k}}{\beta + x_k} - \phi_2(\mu+F) y_{k} + \phi_2\tau_2y_k}{1 + \tau_2\phi_2} = \dfrac{y_k + \phi_2\dfrac{\alpha x_{k}}{\beta + x_k} - \phi_2(\mu+F - \tau_2)y_k}{1 + \tau_2\phi_2}.
%%
\end{split}
\end{equation}
%%
%%
%%
It follows from \eqref{eq:6} that
\begin{equation*}
\begin{split}
&\dfrac{\alpha}{\beta + x_k} +  \mu - \tau_1 \leq \dfrac{\alpha}{\beta} + \mu - \tau_1 \leq 0,\\
%%
%%
%%
&\mu+F - \tau_2 \leq 0.
\end{split}
\end{equation*}
Thus, we deduce from \eqref{eq:7} that $x_{k + 1}, y_{k + 1} \geq 0$ whenever $x_k, y_k \geq 0$. Hence, by mathematical induction, we conclude that $x_k, y_k \geq 0$ for $k > 0$ whenever $x_0, y_0 \geq 0$. This is desired conclusion. The proof is complete.
\end{proof}

\end{document}
